\documentclass[11pt, a4paper]{article}

% ── Paquetes ──────────────────────────────────────────────────────────────────
\usepackage[utf8]{inputenc}
\usepackage[T1]{fontenc}
\usepackage[spanish]{babel}
\usepackage{geometry}
\geometry{margin=2.5cm}
\usepackage{graphicx}
\usepackage{xcolor}
\usepackage{enumitem}
\usepackage{booktabs}
\usepackage{tabularx}
\usepackage{listings}
\usepackage{amsmath}
\usepackage{hyperref}
\usepackage{fancyhdr}
\usepackage{titlesec}
\usepackage{tcolorbox}
\usepackage{fontawesome5}
\usepackage{tikz}
\usetikzlibrary{arrows.meta, positioning, shapes.geometric}

% ── Colores ───────────────────────────────────────────────────────────────────
\definecolor{primary}{HTML}{1E3A5F}
\definecolor{accent}{HTML}{E67E22}
\definecolor{success}{HTML}{27AE60}
\definecolor{danger}{HTML}{E74C3C}
\definecolor{warning}{HTML}{F39C12}
\definecolor{info}{HTML}{3498DB}
\definecolor{codebg}{HTML}{F4F6F9}
\definecolor{codeframe}{HTML}{D5D8DC}

% ── Configuración de listings (JSON) ─────────────────────────────────────────
\lstdefinelanguage{json}{
  basicstyle=\ttfamily\small,
  backgroundcolor=\color{codebg},
  frame=single,
  rulecolor=\color{codeframe},
  breaklines=true,
  showstringspaces=false,
  extendedchars=true,
  inputencoding=utf8,
  literate=
    *{:}{{{\color{accent}{:}}}}{1}
     {,}{{{\color{accent}{,}}}}{1}
     {\{}{{{\color{primary}{\{}}}}{1}
     {\}}{{{\color{primary}{\}}}}}{1}
     {[}{{{\color{primary}{[}}}}{1}
     {]}{{{\color{primary}{]}}}}{1}
     {á}{{{\'a}}}{1}
     {é}{{{\'e}}}{1}
     {í}{{{\'i}}}{1}
     {ó}{{{\'o}}}{1}
     {ú}{{{\'u}}}{1}
     {ñ}{{\~n}}{1}
     {Á}{{{\'A}}}{1}
     {É}{{{\'E}}}{1}
     {Í}{{{\'I}}}{1}
     {Ó}{{{\'O}}}{1}
     {Ú}{{{\'U}}}{1}
     {Ñ}{{\~N}}{1}
}

% ── Configuración de listings (HTTP) ──────────────────────────────────────────
\lstdefinelanguage{http}{
  basicstyle=\ttfamily\small,
  backgroundcolor=\color{codebg},
  frame=single,
  rulecolor=\color{codeframe},
  breaklines=true,
  showstringspaces=false,
  morekeywords={GET,POST,PATCH,DELETE,PUT},
  keywordstyle=\bfseries\color{success},
}

% ── Cajas personalizadas ──────────────────────────────────────────────────────
\newtcolorbox{alertbox}[1][]{
  colback=danger!5, colframe=danger!80!black,
  fonttitle=\bfseries, title={\faExclamationTriangle\quad Corrección Importante}, #1
}

\newtcolorbox{infobox}[1][]{
  colback=info!5, colframe=info!80!black,
  fonttitle=\bfseries, title={\faInfoCircle\quad Nota}, #1
}

\newtcolorbox{tipbox}[1][]{
  colback=success!5, colframe=success!80!black,
  fonttitle=\bfseries, title={\faLightbulb\quad Consejo}, #1
}

\newtcolorbox{warningbox}[1][]{
  colback=warning!5, colframe=warning!80!black,
  fonttitle=\bfseries, title={\faExclamationCircle\quad Atención}, #1
}

% ── Encabezado y pie de página ────────────────────────────────────────────────
\pagestyle{fancy}
\fancyhf{}
\fancyhead[L]{\small\textcolor{primary}{\textbf{Grupo 6} --- Despacho y Logística}}
\fancyhead[R]{\small\textcolor{primary}{Marketplace Cloud}}
\fancyfoot[C]{\small\textcolor{gray}{\thepage}}
\renewcommand{\headrulewidth}{0.4pt}

% ── Formato de títulos ────────────────────────────────────────────────────────
\titleformat{\section}
  {\Large\bfseries\color{primary}}{\thesection.}{0.5em}{}[\titlerule]
\titleformat{\subsection}
  {\large\bfseries\color{primary!80}}{\thesubsection}{0.5em}{}
\titleformat{\subsubsection}
  {\normalsize\bfseries\color{primary!60}}{\thesubsubsection}{0.5em}{}

% ── Hipervínculos ─────────────────────────────────────────────────────────────
\hypersetup{
  colorlinks=true,
  linkcolor=primary,
  urlcolor=info,
  citecolor=accent,
}

% ══════════════════════════════════════════════════════════════════════════════
\begin{document}

% ── Portada ───────────────────────────────────────────────────────────────────
\begin{titlepage}
\centering
\vspace*{3cm}

{\Huge\bfseries\textcolor{primary}{Grupo 6}}\\[0.4cm]
{\LARGE\textcolor{primary!70}{Servicio de Despacho y Logística}}\\[1.5cm]

{\Large Briefing Técnico para el Equipo}\\[0.3cm]
{\large\textcolor{gray}{Documento de alineación previo al desarrollo}}\\[3cm]

\begin{tabular}{rl}
  \textbf{Proyecto:}  & Marketplace en la Nube \\
  \textbf{Asignatura:}& Arquitectura de Software \\
  \textbf{Fecha:}     & 11 de Junio de 2026 \\
  \textbf{Versión:}   & 1.0 --- Revisada y corregida \\
\end{tabular}

\vfill

\begin{warningbox}[title={\faExclamationCircle\quad Antes de leer}]
  Este documento corrige inconsistencias detectadas en el briefing original
  y agrega elementos faltantes. \textbf{Léanlo completo antes de escribir
  una sola línea de código.}
\end{warningbox}

\end{titlepage}

% ── Índice ────────────────────────────────────────────────────────────────────
\tableofcontents
\newpage

% ══════════════════════════════════════════════════════════════════════════════
\section{Contexto del Negocio}
% ══════════════════════════════════════════════════════════════════════════════

Nuestro servicio es el equivalente interno a una empresa de courier (Starken, Chilexpress, Blue Express) dentro del ecosistema del Marketplace. \textbf{No vendemos productos ni procesamos pagos.} Nuestra responsabilidad se limita a tres acciones:

\begin{enumerate}[label=\textbf{\arabic*.}, leftmargin=2em]
  \item \textbf{Crear un despacho:} Cuando un cliente compra exitosamente, el sistema de pedidos nos avisa para que generemos una orden de envío (\textit{Shipment}).
  \item \textbf{Hacer seguimiento:} Mantener en nuestra base de datos el estado actual de cada paquete.
  \item \textbf{Notificar cambios:} Avisar al resto del ecosistema cuando ocurran transiciones de estado.
\end{enumerate}


% ══════════════════════════════════════════════════════════════════════════════
\section{Grupos con los que Interactuamos}
% ══════════════════════════════════════════════════════════════════════════════

\begin{alertbox}[title={\faExclamationTriangle\quad Corrección: Falta G1 en el briefing original}]
  El briefing original solo mencionaba G5, G9 y G10 como aliados, pero en la
  descripción de endpoints se indica que \textbf{el Frontend (G1)} también
  consumirá nuestro \texttt{GET} para mostrar el tracking al usuario.
  Debemos incluirlos.
\end{alertbox}

\begin{table}[h!]
\centering
\renewcommand{\arraystretch}{1.4}
\begin{tabularx}{\textwidth}{c l X l}
  \toprule
  \textbf{Grupo} & \textbf{Nombre} & \textbf{Relación con nosotros} & \textbf{Dirección} \\
  \midrule
  G1  & Frontend       & Consume \texttt{GET /shipments/\{id\}} para mostrar tracking al usuario final. & Ellos $\rightarrow$ Nosotros \\
  G5  & Pedidos        & Llama a \texttt{POST /shipments} cuando un pedido se confirma y paga. Es nuestro \textbf{principal cliente}. & Ellos $\rightarrow$ Nosotros \\
  G9  & Notificaciones & Escucha nuestros eventos (\texttt{ShipmentCreated}, \texttt{ShipmentInTransit}, \texttt{ShipmentDelivered}, etc.) para enviar correos/notificaciones al usuario. & Nosotros $\rightarrow$ Ellos \\
  G10 & Reportería     & Consume nuestros datos (vía eventos o API) para dashboards de rendimiento: entregas a tiempo, tasa de fallos, etc. & Nosotros $\rightarrow$ Ellos \\
  \bottomrule
\end{tabularx}
\caption{Matriz de interacciones del Grupo 6.}
\end{table}

\begin{tipbox}
  \textbf{Acción inmediata:} Contactar a los coordinadores de G5 y G9 para
  acordar los contratos JSON. Propuesta de mensaje:\\[0.3em]
  \textit{``Hola, somos del Grupo 6 (Despachos). Necesitamos definir juntos:
  (1) qué datos nos envían al crear un pedido, y (2) qué datos necesitan
  recibir en las notificaciones de cambio de estado.''}
\end{tipbox}


% ══════════════════════════════════════════════════════════════════════════════
\section{Ciclo de Vida del Envío (Estados)}
\label{sec:estados}
% ══════════════════════════════════════════════════════════════════════════════

\begin{alertbox}[title={\faExclamationTriangle\quad Corrección: Estados faltantes}]
  El briefing original solo contemplaba tres estados lineales:
  \texttt{PENDING} $\rightarrow$ \texttt{IN\_TRANSIT} $\rightarrow$ \texttt{DELIVERED}.
  Esto es insuficiente para cubrir escenarios de error y la rúbrica de
  \textbf{Manejo de Errores}. Se agregan los estados \texttt{CANCELLED},
  \texttt{FAILED} y \texttt{RETURNED}.
\end{alertbox}

\vspace{0.5cm}

% ── Diagrama de estados con TikZ ─────────────────────────────────────────────
\begin{center}
\begin{tikzpicture}[
  node distance=2.8cm,
  state/.style={
    rectangle, rounded corners=6pt, draw=primary, thick,
    fill=primary!8, minimum width=2.8cm, minimum height=1cm,
    font=\small\bfseries, text=primary
  },
  error/.style={
    rectangle, rounded corners=6pt, draw=danger, thick,
    fill=danger!8, minimum width=2.8cm, minimum height=1cm,
    font=\small\bfseries, text=danger
  },
  arrow/.style={-{Stealth[length=3mm]}, thick, primary!70},
  errarrow/.style={-{Stealth[length=3mm]}, thick, danger!70, dashed},
]
  \node[state]                        (pending)    {PENDING};
  \node[state, right of=pending]      (transit)    {IN\_TRANSIT};
  \node[state, right of=transit]      (delivered)  {DELIVERED};
  \node[error, below of=pending]      (cancelled)  {CANCELLED};
  \node[error, below of=transit]      (failed)     {FAILED};
  \node[error, below of=delivered]    (returned)   {RETURNED};

  \draw[arrow]    (pending)   -- (transit);
  \draw[arrow]    (transit)   -- (delivered);
  \draw[errarrow] (pending)   -- (cancelled);
  \draw[errarrow] (transit)   -- (failed);
  \draw[errarrow] (delivered) -- (returned);
\end{tikzpicture}
\end{center}

\vspace{0.3cm}

\begin{table}[h!]
\centering
\renewcommand{\arraystretch}{1.3}
\begin{tabularx}{\textwidth}{l l X}
  \toprule
  \textbf{Estado} & \textbf{Tipo} & \textbf{Descripción} \\
  \midrule
  \texttt{PENDING}     & Normal & Despacho creado, esperando recolección. \\
  \texttt{IN\_TRANSIT} & Normal & Paquete en camino hacia la dirección de destino. \\
  \texttt{DELIVERED}   & Final  & Paquete entregado exitosamente al cliente. \\
  \texttt{CANCELLED}   & Error  & Despacho cancelado antes de salir a reparto (ej: pedido anulado por G5). \\
  \texttt{FAILED}      & Error  & Intento de entrega fallido (dirección incorrecta, destinatario ausente). \\
  \texttt{RETURNED}    & Error  & Paquete devuelto al origen tras entrega fallida o solicitud del cliente. \\
  \bottomrule
\end{tabularx}
\caption{Estados del ciclo de vida de un envío.}
\end{table}


% ══════════════════════════════════════════════════════════════════════════════
\section{Contrato REST (API de Shipments)}
\label{sec:api}
% ══════════════════════════════════════════════════════════════════════════════

\begin{alertbox}[title={\faExclamationTriangle\quad Corrección: Endpoint PATCH faltante}]
  El briefing original solo definía \texttt{POST} y \texttt{GET}. Sin embargo,
  para actualizar el estado del envío (por ejemplo, cuando la simulación
  de la Fase 4 avanza el paquete de \texttt{IN\_TRANSIT} a \texttt{DELIVERED}),
  se necesita un endpoint \texttt{PATCH}. Se agrega a continuación.
\end{alertbox}

\subsection{Base URL}
\begin{lstlisting}[language=http]
https://<nombre-del-servicio>.onrender.com/api/v1
\end{lstlisting}

% ── POST ──────────────────────────────────────────────────────────────────────
\subsection{POST /shipments --- Crear un despacho}

\textbf{Quién lo llama:} Grupo 5 (Pedidos).

\subsubsection*{Request Body}
\begin{lstlisting}[language=json]
{
  "order_id": "ORD-20260611-001",
  "customer_name": "María González",
  "address": "Av. Providencia 1234, Depto 56",
  "city": "Santiago",
  "weight_kg": 3.2
}
\end{lstlisting}

\subsubsection*{Response --- 201 Created}
\begin{lstlisting}[language=json]
{
  "shipment_id": "SHP-a1b2c3d4",
  "order_id": "ORD-20260611-001",
  "status": "PENDING",
  "created_at": "2026-06-11T14:00:00Z",
  "estimated_delivery": "2026-06-14T18:00:00Z"
}
\end{lstlisting}

\subsubsection*{Errores posibles}
\begin{table}[h!]
\centering
\begin{tabularx}{\textwidth}{l l X}
  \toprule
  \textbf{Código} & \textbf{Nombre} & \textbf{Cuándo} \\
  \midrule
  400 & Bad Request             & Faltan campos obligatorios o el JSON es inválido. \\
  409 & Conflict                & Ya existe un shipment para ese \texttt{order\_id}. \\
  422 & Unprocessable Entity    & Datos presentes pero inválidos (ej: \texttt{weight\_kg} negativo). \\
  \bottomrule
\end{tabularx}
\end{table}

% ── GET ───────────────────────────────────────────────────────────────────────
\subsection{GET /shipments/\{shipment\_id\} --- Consultar estado}

\textbf{Quién lo llama:} Grupo 1 (Frontend) y Grupo 5 (Pedidos).

\subsubsection*{Response --- 200 OK}
\begin{lstlisting}[language=json]
{
  "shipment_id": "SHP-a1b2c3d4",
  "order_id": "ORD-20260611-001",
  "status": "IN_TRANSIT",
  "customer_name": "Maria Gonzalez",
  "address": "Av. Providencia 1234, Depto 56",
  "city": "Santiago",
  "weight_kg": 3.2,
  "created_at": "2026-06-11T14:00:00Z",
  "updated_at": "2026-06-11T14:30:00Z",
  "estimated_delivery": "2026-06-14T18:00:00Z"
}
\end{lstlisting}

\subsubsection*{Errores posibles}
\begin{table}[h!]
\centering
\begin{tabularx}{\textwidth}{l l X}
  \toprule
  \textbf{Código} & \textbf{Nombre} & \textbf{Cuándo} \\
  \midrule
  404 & Not Found & No existe un shipment con ese ID. \\
  \bottomrule
\end{tabularx}
\end{table}

% ── PATCH ─────────────────────────────────────────────────────────────────────
\subsection{PATCH /shipments/\{shipment\_id\} --- Actualizar estado}

\textbf{Quién lo llama:} Nuestro propio servicio (simulación automática de la Fase 4) o un endpoint administrativo.

\subsubsection*{Request Body}
\begin{lstlisting}[language=json]
{
  "status": "DELIVERED"
}
\end{lstlisting}

\subsubsection*{Response --- 200 OK}
\begin{lstlisting}[language=json]
{
  "shipment_id": "SHP-a1b2c3d4",
  "status": "DELIVERED",
  "updated_at": "2026-06-11T15:00:00Z"
}
\end{lstlisting}

\subsubsection*{Errores posibles}
\begin{table}[h!]
\centering
\begin{tabularx}{\textwidth}{l l X}
  \toprule
  \textbf{Código} & \textbf{Nombre} & \textbf{Cuándo} \\
  \midrule
  400 & Bad Request             & Estado enviado no es válido. \\
  404 & Not Found               & No existe un shipment con ese ID. \\
  409 & Conflict                & Transición de estado inválida (ej: de \texttt{DELIVERED} a \texttt{PENDING}). \\
  \bottomrule
\end{tabularx}
\end{table}

\begin{infobox}
  Las transiciones válidas son las mostradas en el diagrama de la
  Sección~\ref{sec:estados}. Cualquier otra transición debe rechazarse
  con un \texttt{409 Conflict}.
\end{infobox}


% ── G6 / G9 / G10 Eventos ─────────────────────────────────────────────────────
\section{Contrato de Eventos}
\label{sec:eventos}
% ══════════════════════════════════════════════════════════════════════════════

\begin{alertbox}[title={\faExclamationTriangle\quad Corrección: Mecanismo de eventos no especificado}]
  El briefing original habla de ``publicar eventos'' y ``emitir mensajes''
  sin aclarar la tecnología. Para un proyecto desplegado en Render (tier
  gratuito), las opciones realistas son:
  \begin{itemize}[noitemsep]
    \item \textbf{Webhooks HTTP (recomendado):} Nuestro servicio hace un
          \texttt{POST} a una URL registrada del G9 cada vez que cambia un estado.
    \item \textbf{Polling:} G9/G10 consultan periódicamente nuestro \texttt{GET}.
          Menos eficiente pero más simple.
  \end{itemize}
  \textbf{Decisión del equipo:} Debemos acordar esto con G9 antes del 16 de Junio.
\end{alertbox}

\subsection{Formato estándar del evento}

Cada evento que publiquemos debe seguir esta estructura:

\begin{lstlisting}[language=json]
{
  "event_type": "ShipmentDelivered",
  "event_id": "evt-uuid-autogenerado",
  "timestamp": "2026-06-11T15:00:00Z",
  "payload": {
    "shipment_id": "SHP-a1b2c3d4",
    "order_id": "ORD-20260611-001",
    "previous_status": "IN_TRANSIT",
    "new_status": "DELIVERED",
    "customer_name": "Maria Gonzalez",
    "city": "Santiago"
  }
}
\end{lstlisting}

\subsection{Catálogo de eventos}

\begin{table}[h!]
\centering
\renewcommand{\arraystretch}{1.3}
\begin{tabularx}{\textwidth}{l l X}
  \toprule
  \textbf{event\_type} & \textbf{Transición} & \textbf{Consumidor principal} \\
  \midrule
  \texttt{ShipmentCreated}    & $\emptyset \rightarrow$ \texttt{PENDING}     & G9 (notifica al cliente), G10 \\
  \texttt{ShipmentInTransit}  & \texttt{PENDING} $\rightarrow$ \texttt{IN\_TRANSIT}  & G9 (``Tu paquete va en camino'') \\
  \texttt{ShipmentDelivered}  & \texttt{IN\_TRANSIT} $\rightarrow$ \texttt{DELIVERED} & G9 (``Tu paquete llegó''), G10 \\
  \texttt{ShipmentCancelled}  & \texttt{PENDING} $\rightarrow$ \texttt{CANCELLED}    & G9, G5 (confirmar cancelación) \\
  \texttt{ShipmentFailed}     & \texttt{IN\_TRANSIT} $\rightarrow$ \texttt{FAILED}    & G9 (``Hubo un problema''), G5 \\
  \texttt{ShipmentReturned}   & \texttt{DELIVERED} $\rightarrow$ \texttt{RETURNED}    & G9, G5, G10 \\
  \bottomrule
\end{tabularx}
\caption{Catálogo completo de eventos del servicio de despacho.}
\end{table}


% ══════════════════════════════════════════════════════════════════════════════
\section{Conceptos Técnicos Obligatorios}
% ══════════════════════════════════════════════════════════════════════════════

\begin{table}[h!]
\centering
\renewcommand{\arraystretch}{1.4}
\begin{tabularx}{\textwidth}{l X}
  \toprule
  \textbf{Concepto} & \textbf{Cómo lo demostramos} \\
  \midrule
  \textbf{REST} &
    Exponemos los endpoints \texttt{POST}, \texttt{GET} y \texttt{PATCH} descritos
    en la Sección~\ref{sec:api}, siguiendo convenciones REST (verbos HTTP, códigos
    de estado, recursos en plural). \\
  \textbf{Eventos} &
    Emitimos eventos (Sección~\ref{sec:eventos}) en cada transición de estado,
    usando webhooks HTTP o el mecanismo acordado con G9. \\
  \textbf{Consistencia Eventual} &
    El despacho no es instantáneo. El usuario ve \texttt{PENDING} inmediatamente,
    y los estados cambian de forma asíncrona (simulación temporizada). Los otros
    grupos se enteran vía eventos, no por consulta síncrona. \\
  \textbf{Manejo de Errores} &
    Validamos inputs, respondemos con códigos HTTP correctos (400, 404, 409, 422),
    y contemplamos estados de error (\texttt{CANCELLED}, \texttt{FAILED},
    \texttt{RETURNED}). \\
  \bottomrule
\end{tabularx}
\end{table}


% ══════════════════════════════════════════════════════════════════════════════
\section{Hoja de Ruta por Fases}
% ══════════════════════════════════════════════════════════════════════════════

\begin{alertbox}[title={\faExclamationTriangle\quad Corrección crítica: Fecha de Fase 3}]
  El briefing original indicaba que la Fase 3 vence el \textbf{30 de Julio},
  pero la Fase 4 vence el \textbf{7 de Julio}. Esto es cronológicamente
  imposible. La fecha correcta de la Fase 3 es \textbf{30 de Junio}.
\end{alertbox}

\begin{table}[h!]
\centering
\renewcommand{\arraystretch}{1.5}
\begin{tabularx}{\textwidth}{c l c X}
  \toprule
  \textbf{Fase} & \textbf{Nombre} & \textbf{Vencimiento} & \textbf{Entregables} \\
  \midrule
  1 & Contratos &
    \textbf{16 Jun} &
    Contrato REST (Swagger/Markdown), contrato de eventos (JSON), diagrama de
    arquitectura, matriz de dependencias. \\
  2 & Mock Funcional &
    \textbf{23 Jun} &
    Repositorio en GitHub. Servicio desplegado en Render que responde JSON
    simulado a los endpoints, sin lógica real. Permite al G5 probar su código. \\
  3 & Despliegue Real &
    \textcolor{danger}{\textbf{30 Jun}} &
    Servicio funcional (FastAPI o Node.js) en Render, conectado a base de datos
    en Supabase/Firebase. Tabla \texttt{shipments} con persistencia real. \\
  4 & Integración &
    \textbf{7 Jul} &
    Conexión real con G5 y G9. Script de simulación temporal:
    \texttt{PENDING} $\xrightarrow{\text{30s}}$ \texttt{IN\_TRANSIT}
    $\xrightarrow{\text{60s}}$ \texttt{DELIVERED}, emitiendo eventos en cada paso. \\
  5 & Defensa Final &
    \textbf{13 Jul} &
    Presentación con diapositivas, demo en vivo del flujo completo
    (pedido $\rightarrow$ despacho $\rightarrow$ entrega), defensa oral
    individual (2 preguntas técnicas por persona). \\
  \bottomrule
\end{tabularx}
\caption{Cronograma corregido del proyecto.}
\end{table}


% ══════════════════════════════════════════════════════════════════════════════
\section{Stack Tecnológico Propuesto}
% ══════════════════════════════════════════════════════════════════════════════

\begin{table}[h!]
\centering
\renewcommand{\arraystretch}{1.3}
\begin{tabularx}{\textwidth}{l l X}
  \toprule
  \textbf{Capa} & \textbf{Tecnología} & \textbf{Justificación} \\
  \midrule
  Framework    & Python + FastAPI  & Rápido de desarrollar, tipado, documentación automática (Swagger). \\
  Hosting      & Render (Free)     & Tier gratuito suficiente para el proyecto. \\
  Base de datos& Supabase (PostgreSQL) & Tier gratuito, interfaz web para inspección, compatible con SQL estándar. \\
  Eventos      & Webhooks HTTP     & Sin infraestructura adicional; un simple \texttt{POST} a la URL de G9. \\
  Versionado   & GitHub            & Repositorio del equipo con ramas por fase. \\
  \bottomrule
\end{tabularx}
\end{table}


% ══════════════════════════════════════════════════════════════════════════════
\section{Esquema de Base de Datos}
% ══════════════════════════════════════════════════════════════════════════════

\begin{lstlisting}[language=json, title={\textbf{Tabla: shipments}}]
{
  "shipment_id":        "TEXT PRIMARY KEY  (ej: SHP-uuid)",
  "order_id":           "TEXT UNIQUE NOT NULL",
  "customer_name":      "TEXT NOT NULL",
  "address":            "TEXT NOT NULL",
  "city":               "TEXT NOT NULL",
  "weight_kg":          "REAL NOT NULL CHECK (weight_kg > 0)",
  "status":             "TEXT NOT NULL DEFAULT 'PENDING'",
  "created_at":         "TIMESTAMP DEFAULT NOW()",
  "updated_at":         "TIMESTAMP DEFAULT NOW()",
  "estimated_delivery": "TIMESTAMP"
}
\end{lstlisting}

\begin{infobox}
  Usar un \texttt{CHECK} constraint en \texttt{status} para que solo acepte
  los valores válidos: \texttt{PENDING}, \texttt{IN\_TRANSIT},
  \texttt{DELIVERED}, \texttt{CANCELLED}, \texttt{FAILED}, \texttt{RETURNED}.
\end{infobox}


% ══════════════════════════════════════════════════════════════════════════════
\section{Diagrama de Arquitectura}
% ══════════════════════════════════════════════════════════════════════════════

\begin{center}
\begin{tikzpicture}[
  node distance=2.5cm and 3.5cm,
  box/.style={
    rectangle, rounded corners=8pt, draw=primary, thick, fill=primary!6,
    minimum width=3cm, minimum height=1.4cm, align=center,
    font=\small\bfseries, text=primary
  },
  extbox/.style={
    rectangle, rounded corners=8pt, draw=info, thick, fill=info!6,
    minimum width=3cm, minimum height=1.4cm, align=center,
    font=\small\bfseries, text=info
  },
  db/.style={
    cylinder, draw=accent, thick, fill=accent!8, shape border rotate=90,
    minimum width=2cm, minimum height=1.2cm, aspect=0.3,
    font=\small\bfseries, text=accent
  },
  arr/.style={-{Stealth[length=3mm]}, thick},
]
  % Nodos
  \node[extbox]                          (g1)  {G1\\\relax Frontend};
  \node[extbox, below=1.5cm of g1]       (g5)  {G5\\\relax Pedidos};
  \node[box, right=3cm of g5]            (g6)  {G6\\\relax Despachos};
  \node[db, below=1.5cm of g6]           (db)  {Supabase};
  \node[extbox, right=3cm of g6]         (g9)  {G9\\\relax Notificaciones};
  \node[extbox, below=1.5cm of g9]       (g10) {G10\\\relax Reportería};

  % Flechas
  \draw[arr, info] (g5)  -- node[above, font=\scriptsize] {\texttt{POST}} (g6);
  \draw[arr, info] (g1)  -| node[above left, font=\scriptsize] {\texttt{GET}} (g6);
  \draw[arr, primary] (g6) -- node[left, font=\scriptsize] {SQL} (db);
  \draw[arr, success] (g6) -- node[above, font=\scriptsize] {Webhook} (g9);
  \draw[arr, success] (g6) |- node[below right, font=\scriptsize] {Evento} (g10);
\end{tikzpicture}
\end{center}


% ══════════════════════════════════════════════════════════════════════════════
\section{Resumen de Correcciones Aplicadas}
% ══════════════════════════════════════════════════════════════════════════════

A continuación se listan todas las correcciones y adiciones respecto al
briefing original:

\begin{enumerate}[label=\textcolor{danger}{\textbf{C\arabic*.}}, leftmargin=2.5em]
  \item \textbf{Fecha de Fase 3:} Cambiada de ``30 de Julio'' a
        ``\textbf{30 de Junio}''. La fecha original hacía que la Fase 4
        (7 de Julio) venciera antes que la Fase 3.

  \item \textbf{Grupo 1 como consumidor:} Agregado G1 (Frontend) a la
        matriz de interacciones, ya que consume el endpoint
        \texttt{GET /shipments/\{id\}} para tracking.

  \item \textbf{Estados de error:} Agregados \texttt{CANCELLED},
        \texttt{FAILED} y \texttt{RETURNED} al ciclo de vida del envío.
        El briefing original solo contemplaba el ``camino feliz''.

  \item \textbf{Endpoint PATCH:} Agregado
        \texttt{PATCH /shipments/\{id\}} para actualizar el estado del
        envío. Esencial para la simulación automática de la Fase 4.

  \item \textbf{Mecanismo de eventos:} Especificado que usaremos
        \textbf{webhooks HTTP} como mecanismo por defecto. El briefing
        original no aclaraba la tecnología.

  \item \textbf{Payloads de eventos:} Agregados ejemplos concretos del
        formato JSON de los eventos, con campos \texttt{event\_type},
        \texttt{event\_id}, \texttt{timestamp} y \texttt{payload}.

  \item \textbf{Esquema de base de datos:} Agregada la definición de la
        tabla \texttt{shipments} con tipos, constraints y valores por
        defecto.

  \item \textbf{Códigos de error detallados:} Documentados los errores
        HTTP esperados para cada endpoint (400, 404, 409, 422).
\end{enumerate}


% ══════════════════════════════════════════════════════════════════════════════
\section{Próximos Pasos Inmediatos}
% ══════════════════════════════════════════════════════════════════════════════

\begin{enumerate}[label=\textbf{\arabic*.}]
  \item \textbf{Hoy (11 Jun):} Cada miembro del equipo lee este documento
        completo.
  \item \textbf{Hoy/Mañana:} Contactar a G5 y G9 para acordar contratos
        JSON.
  \item \textbf{Antes del 14 Jun:} Generar el documento Swagger/Markdown
        definitivo con los contratos REST y de eventos.
  \item \textbf{Antes del 16 Jun:} Entregar Fase 1 (contratos + diagrama
        de arquitectura + matriz de dependencias).
  \item \textbf{17--23 Jun:} Crear repo GitHub, implementar el mock en
        Render y entregar Fase 2.
\end{enumerate}

\begin{tipbox}
  Si tienen dudas sobre algún punto de este documento, discútanlo en el
  grupo \textbf{antes} de asumir cosas. Las decisiones de diseño que
  tomemos ahora nos ahorrarán reescrituras después.
\end{tipbox}

\end{document}
